\section{事前登録の経緯}\label{app:prereg}
判定の基準は，計算や学習の前に実験の記録に書いて保存した．
本文の各指標の位置づけを表\ref{tab:metrics}に示し，ここでは結果を見た後に変えたことと，その理由をまとめる．

\paragraph{他者側の主の指標}
段階Iの他者側の関門として，G5a（保った快・不快のもとでの他者側の傾き）と G5b（$\pm a$ の差）を事前登録した．
経験50\%条件では両方に合格したが，経験80\%条件では G5b に合格し，G5a は3組中1組で不合格であった．
経験80\%条件の前に登録した予測「G5a の傾きが経験50\%条件より大きくなる」も外れた．
その後の事前登録外の分析で，G5a の傾きには快・不快の符号に従う動きと，保つ前に置いた位置から本当の位置へ戻る動き（符号に依らない動き）が混ざることが分かった．
経験50\%条件では戻る動きが傾きの半分ほどを占めて押し上げ，経験80\%条件では逆向きの戻る動きが符号に従う動きを打ち消していたと読める（推測）．
そこで，G5a を副に下げ，同じ考え方で戻る動きを除いた「符号に従う分」を定義した．
この判断は結果を見た後のものであり，符号に従う分を事前登録してから計算したのは組5・6 だけである．
G5b は，同じ状態からの $+a$ と $-a$ の差をとるので戻る動きの影響を受けず，当初から事前登録した主の関門であったため，本文では他者側の主の根拠とした．
経験50\%条件の G5a の合格は，好みに従う動きの根拠には用いない．
% 出典：v13 まとめ §3.1，執筆メモ §4 の H2b の注，v13 まとめ §3 表1。

\paragraph{組の追加と組をまたいだ統計}
判定は本番の3組で行い，予備の組4と，後から追加した組5・6 は，追加の前に登録したとおり除外せずに報告した．
組をまたいだ差の $t$ 分布による信頼区間は，事前登録していない記述的な分析である．
組が6つでは，すべての組で正でも符号検定の $p$ は 0.031 にしかならない．
% 出典：v13 まとめ §10・§11，v14 ログ §43。

\paragraph{段階IIの扱い}
事前登録した規則では，段階Iの関門のうち G5a に不合格だった条件では，段階IIを探索として行う．
経験80\%条件の自由出力モジュールはこれに当たる．
形制限モジュールは自由出力モジュールの結果を見た後に設計した探索であり，他者遮蔽学習は三層の分析から導いた仮説を事前登録して確かめた探索である．
1ビットモジュールは事前登録した．
ただし経験50\%条件の上での方角との比較は，判定を見た後に方法を決めた確認である．
% 出典：v13 まとめ §4，v12 まとめ §3，v14 まとめ §6。

\paragraph{学習後に重みを写す比較}
非共有条件のモデルに学習後に自分側の重みを写す比較は，計算の前に次を登録した．
予想は，(1) すべて写したモデルの G5b が元の非共有モデルより大きく，経験なし条件との差が0を含まない（本番の3組中2組以上で支持），(2) 写したモデルの G5b・符号に従う分・照合の BA は経験80\%条件に近づく（記述），である．
あわせて，写すと他者の追跡が壊れうるので，$h^2$ から A-2 の位置を読む決定係数が 0.8 を下回る組は G5b を判定しないと決めた．
快・不快の列だけを写す条件は，結果を見る前に追加で登録した．
経験なし条件のモデルに写す確認は事前登録外である．
% 出典：docs/paper_extra_log.md §3・§13・§18。

\paragraph{仮説の照合}
照合の手続き，候補，主の指標（自己視界基準の BA，$T=50$，L1 期間の最終5点）と判定は，計算の前に登録した．
他者特徴基準の数え方（MSE 期間の最終5点）と，方角との食い違いによる好み追従基準は，経験50\%条件の結果を見た後，経験80\%条件の結果を見る前に登録した．
他者特徴基準の値は L1 期間に変わらないので，当初の数え方では同じ値を5回数えていたことが分かり，訂正した．
% 出典：v14 まとめ §6，v14 ログ §20・§21。

\paragraph{投影の誤り}
真値条件の悪化は予想していなかった結果であり，その解釈に用いた分析（一歩の種類ごとの差，保つ前の A-2 の置き方）は事前登録外である．
% 出典：v13 まとめ §4，§3.1。

\paragraph{他者の行動の仕方を変えた比較}
真値条件の悪化を見た後に，その原因を調べるために計画した比較であり，計算の前に次を登録した．
主の量は真値条件の誤差から零を入れた誤差を引いた $\Delta$ である．
予想は三つで，(1) 経験80\%条件では，A-2 の規則の確率を 0.5 から 0.8 にすると $\Delta$ が小さくなる，(2) その変化は経験50\%条件より大きい（交互作用），(3) 経験なし条件との交互作用も同じ向きである．
判定は，各組で同じエピソードを対にしたブートストラップの 95\% 信頼区間が0を含まず予想の向きである組が本番の3組中2組以上なら支持とし，緑好みデータと好み混合データで別々に行い，片方だけで成り立てば一部支持とした．
全6組では組をまたいだ $t$ の信頼区間と成り立つ組の数を並べた．
結果は本文4.5節に示した（三つとも両データで支持）．
% 出典：docs/paper_extra_log.md §2（事前登録，Claude Code）。
